% ==============================================================================
% PDF-1 / Section 07 — Non-Claims and Limits (Skeleton)
% No new primitives. Diagnostic-only. Container structure.
% ==============================================================================

\section{Non-Claims and Limits}
\label{sec:nonclaims-limits}

\subsection{Explicit non-claims}
This document makes the following explicit non-claims:
\begin{itemize}
  \item No prescriptive guidance for enterprise change or intervention.
  \item No optimization objectives or target states.
  \item No performance benchmarks, KPIs, or maturity models.
  \item No predictive guarantees or probabilistic forecasts.
  \item No managerial, governance, or policy recommendations.
\end{itemize}

\subsection{Limits of applicability}
The diagnostic constructions presented here are applicable only within the
formal scope defined by KEA (PDF-0) and \texttt{ETS.K\_EA.v1.1}.
Outside this scope, diagnostic results may be undefined, non-falsifiable,
or collapse into STOP regions.

This section does not attempt to extend applicability beyond formally
specified boundaries.

\subsection{Epistemic status}
All results obtained via this PDF-1 are diagnostic descriptions under
explicit assumptions.
They do not constitute explanations, causes, or actionable knowledge.

Uncertainty, partial observability, and model mismatch are treated as
first-class limiting factors, not as defects to be corrected.

\subsection{Relation to future runs}
Any future extension (e.g., computational demos, executable traces)
must:
\begin{itemize}
  \item preserve all non-claims stated here,
  \item remain compatible with invariant stress tests,
  \item avoid introducing implicit prescriptions.
\end{itemize}

\subsection{Closure}
This section serves as a hard boundary for interpretation.
Crossing this boundary converts the document into a different class of artifact,
which is explicitly out of scope for PDF-1.
